% =================================================
% Ученический рабочий лист. Урок 2
% =================================================

\worksheettitle
  {Классы. Чтение и запись чисел}
  {Урок 2 из 3. Обозначение натуральных чисел}

\studentnameblock

\begin{checkpointbox}[Входная диагностика: 4 минуты]
  \begin{enumerate}[label=\textbf{\arabic*.},itemsep=0.38em]
    \item Какая цифра стоит в разряде тысяч числа \(62\,405\)?
      Ответ: \answerblank[15mm].
    \item \(70\,308=\answerblank[58mm].\)
    \item \(40\,000+700+2=\answerblank[24mm].\)
    \item Сколько цифр в записи числа \(5\,008\,019\)?
      Ответ: \answerblank[15mm].
  \end{enumerate}
\end{checkpointbox}

\section{Делим число на классы}

Разделите записи справа налево на группы по три цифры. Поставьте пробелы:
\[
  5308\longrightarrow\answerblank[35mm],\qquad
  7204061\longrightarrow\answerblank[42mm],
\]
\[
  4805070204\longrightarrow\answerblank[52mm].
\]

\splitclassesfigure

\begin{reflectionbox}[Сформулируйте правило]
  Число делят на классы, начиная с \answerblank[32mm] края,
  по \answerblank[13mm] цифры.
\end{reflectionbox}

\section{Названия классов}

Заполните пропуски:
\[
  \underbrace{\text{единицы}}_{\text{I класс}},\qquad
  \underbrace{\answerblank[24mm]}_{\text{II класс}},\qquad
  \underbrace{\answerblank[28mm]}_{\text{III класс}},\qquad
  \underbrace{\answerblank[30mm]}_{\text{IV класс}}.
\]

\classtablefigure

\begin{fillbox}[Разберите число \(7\,304\,052\,018\)]
  \begin{enumerate}
    \item В классе миллиардов записано число \answerblank[14mm].
    \item В классе миллионов записано число \answerblank[18mm].
    \item В классе тысяч записана группа \answerblank[18mm], то есть число \answerblank[14mm].
    \item В классе единиц записана группа \answerblank[18mm], то есть число \answerblank[14mm].
  \end{enumerate}
\end{fillbox}

\section{Работаем с разрядной таблицей}

Рассмотрите число \(325\,041\,708\).

\classworksheetfigure

\begin{fillbox}[Прочитайте таблицу]
  \begin{enumerate}
    \item Класс миллионов содержит число \answerblank[18mm].
    \item Класс тысяч содержит число \answerblank[18mm].
    \item Класс единиц содержит число \answerblank[18mm].
    \item Цифра сотен тысяч равна \answerblank[12mm].
    \item Цифра десятков единиц равна \answerblank[12mm].
  \end{enumerate}
\end{fillbox}

\begin{thinkbox}[Группа и число внутри класса]
  Почему в таблице класс тысяч записан как \(041\), хотя мы говорим
  «сорок одна тысяча»?
  \writinglines{2}
\end{thinkbox}

\section{Читаем многозначные числа}

\readingroutefigure

Запишите числа словами.

\begin{fillbox}
  \begin{enumerate}
    \item \(4\,218\):
      \writinglines{1}
    \item \(73\,005\):
      \writinglines{1}
    \item \(6\,040\,019\):
      \writinglines{1}
    \item \(502\,007\,030\):
      \writinglines{1}
  \end{enumerate}
\end{fillbox}

\begin{checkpointbox}[Как читать класс \(000\)]
  Рассмотрите число \(9\,000\,406\).
  \zeroclassfigure
  \begin{enumerate}
    \item Какой класс содержит \(000\)? \answerblank[30mm].
    \item Произносим ли мы название этого класса? \choicebox\ да \quad \choicebox\ нет.
    \item Запишите число словами:
      \writinglines{1}
  \end{enumerate}
\end{checkpointbox}

\section{Записываем число цифрами}

\writingclassesfigure

\begin{fillbox}[Сначала заполните классы]
  \renewcommand{\arraystretch}{1.55}
  \begin{tabularx}{\linewidth}{@{}X>{\centering\arraybackslash}p{0.16\linewidth}>{\centering\arraybackslash}p{0.16\linewidth}>{\centering\arraybackslash}p{0.16\linewidth}@{}}
    \toprule
    \rowcolor{TableHeader}
    Число словами & миллионы & тысячи & единицы \\
    \midrule
    сорок два миллиона семь тысяч девяносто
      & \answerblank[15mm] & \answerblank[15mm] & \answerblank[15mm] \\
    триста пять миллионов восемьсот два
      & \answerblank[15mm] & \answerblank[15mm] & \answerblank[15mm] \\
    восемнадцать миллионов сорок тысяч пять
      & \answerblank[15mm] & \answerblank[15mm] & \answerblank[15mm] \\
    \bottomrule
  \end{tabularx}
\end{fillbox}

Теперь запишите числа полностью:
\[
  \answerblank[42mm],\qquad
  \answerblank[42mm],\qquad
  \answerblank[42mm].
\]

\begin{fillbox}[Запишите цифрами]
  \begin{enumerate}
    \item шесть миллионов сорок два: \answerblank[34mm];
    \item семьдесят миллионов три тысячи восемь: \answerblank[38mm];
    \item девять миллиардов три миллиона двадцать тысяч четыре.\par
      Ответ: \answerblank[46mm].
  \end{enumerate}
\end{fillbox}

\section{Исправляем ошибки}

\begin{warningbox}[Ошибка 1]
  Ученик записал:
  \[
    \text{шесть миллионов сорок два}=6\,042.
  \]
  Какой класс потерян? \answerblank[35mm].

  Верная запись: \answerblank[40mm].
\end{warningbox}

\begin{warningbox}[Ошибка 2]
  Ученик прочитал \(53\,008\,004\) так:
  «пятьдесят три миллиона восемьсот четыре».

  Какие слова он пропустил? \answerblank[55mm].

  Верное чтение:
  \writinglines{1}
\end{warningbox}

\section{Самостоятельная проверка}

\begin{practicebox}[Выполните без подсказок]
  \begin{enumerate}[label=\textbf{\arabic*.},itemsep=0.4em]
    \item Разделите на классы: \(805070204\rightarrow\answerblank[45mm]\).
    \item Какое число записано в классе тысяч числа \(71\,006\,540\)?
      \answerblank[18mm].
    \item Прочитайте число \(40\,005\,018\):
      \writinglines{1}
    \item Запишите цифрами: двести три миллиона семьдесят тысяч девять.
      \answerblank[42mm].
    \item Запишите цифрами: четыре миллиарда двенадцать миллионов пять.
      \answerblank[45mm].
  \end{enumerate}
\end{practicebox}

\begin{checkpointbox}[Выходной билет]
  \begin{enumerate}
    \item В числе \(8\,014\,005\) класс тысяч записан группой \answerblank[18mm].
    \item Число \(3\,000\,207\) читается:
      \writinglines{1}
    \item «Пятьдесят миллионов восемь» записывается \answerblank[35mm].
  \end{enumerate}
\end{checkpointbox}

\begin{reflectionbox}[Оцените себя]
  \choicebox\ могу объяснить другому;
  \qquad
  \choicebox\ умею по образцу;
  \qquad
  \choicebox\ нужна тренировка.
\end{reflectionbox}
